% =============================================================================
% SECCIÓN 3: RESULTADOS
% Capítulo 2: Instrumentación IoT y Caracterización Metrológica
% =============================================================================


% -----------------------------------------------------------------------------
% 3.1. CARACTERIZACIÓN FISICOQUÍMICA
% -----------------------------------------------------------------------------
\subsection{Caracterización fisicoquímica de sustratos locales}
\label{subsec:cap2_caracterizacion_sustratos}

La composición proximal de los sustratos agroindustriales empleados
reveló diferencias marcadas entre la dieta D1 (base proteica, residuo
de procesamiento avícola) y la dieta D4 (base fibrosa, mezcla de
residuos vegetales del Valle del Cauca), con implicaciones directas
sobre la cinética de degradación y el perfil de emisiones gaseosas
(Tabla~\ref{tab:composicion_sustratos}).

\begin{table}[htbp]
\centering
\caption{Composición proximal de los sustratos experimentales.
Valores expresados en base húmeda, excepto C:N (relación molar).}
\label{tab:composicion_sustratos}
\begin{tabular}{lcc}
\toprule
\textbf{Parámetro} & \textbf{Dieta D1 (Proteica)} & \textbf{Dieta D4 (Fibrosa)} \\
\midrule
Humedad (\%) & 68.2 $\pm$ 1.4 & 72.5 $\pm$ 2.1 \\
Cenizas (\%) & 8.5 $\pm$ 0.3 & 6.1 $\pm$ 0.4 \\
Proteína bruta (\%) & 22.4 $\pm$ 0.8 & 9.3 $\pm$ 0.5 \\
Lípidos totales (\%) & 12.1 $\pm$ 0.6 & 3.2 $\pm$ 0.3 \\
Fibra bruta (\%) & 5.8 $\pm$ 0.4 & 28.7 $\pm$ 1.2 \\
Carbohidratos solubles (\%) & 14.2 $\pm$ 0.9 & 11.2 $\pm$ 0.7 \\
pH inicial & 6.4 $\pm$ 0.1 & 5.9 $\pm$ 0.2 \\
Relación C:N & 8.2 $\pm$ 0.3 & 21.4 $\pm$ 0.8 \\
\bottomrule
\end{tabular}
\end{table}

La dieta D1 presentó una relación C:N cercana al óptimo reportado para
la maximización de la tasa de conversión en BSF (C:N $\approx$ 7--10),
con alto contenido de proteína y lípidos que favorecen el crecimiento
larval acelerado \autocite{nayakHermetiaIllucensDiptera2024}. En
contraste, la dieta D4 exhibió una relación C:N elevada ($>20$) y alto
contenido de fibra bruta, condiciones asociadas a menores tasas de
asimilación y mayores tiempos de residencia del sustrato
\autocite{siddiquiBlackSoldierFly2022}. El pH inicial ligeramente ácido
de D4 (5.9) contrastó con el valor cercano a la neutralidad de D1 (6.4),
factor que condiciona la volatilización de amoníaco y la actividad
microbiana asociada durante la bioconversión.

% -----------------------------------------------------------------------------
% 3.2. ESTABILIDAD Y DERIVA
% -----------------------------------------------------------------------------
\subsection{Estabilidad y deriva de sensores en HR$>80$\%}
\label{subsec:cap2_deriva_sensores}

La exposición continua del nodo sensor a humedad relativa superior al
80\,\% durante los ciclos de bioconversión indujo una deriva
sistemática positiva en las lecturas de \ce{CO2}, mientras que el
sensor de \ce{CH4} mostró una deriva negativa progresiva tras las
primeras 48 h de operación (Figura~\ref{fig:deriva_sensores}).

\begin{figure}[htbp]
\centering
% \includegraphics[width=0.8\textwidth]{figuras/deriva_sensores_hr80}
\caption{Deriva relativa de sensores NDIR durante exposición continua
a HR $> 80$\,\%. Línea punteada: umbral de aceptabilidad metrológica
($\pm 5$\,\% de fondo de escala).}
\label{fig:deriva_sensores}
\end{figure}

El sensor SCD41 para \ce{CO2} exhibió una deriva del $+3.2$\,\% del
fondo de escala cada 24 h en HR promedio de 82\,\%, valor que superó
la especificación del fabricante para condiciones normales (HR $< 60$\,\%),
pero que permaneció dentro del rango corregible mediante el modelo de
compensación implementado (Ecuación~\ref{eq:compensacion_ndir}). La
deriva fue altamente correlacionada con la humedad relativa ($R = 0.91$,
$p < 0.001$) y débilmente con la temperatura ($R = 0.34$), confirmando
que la humedad es el factor dominante de perturbación metrológica en
este entorno. El tiempo de estabilización tras cada evento de
condensación transitoria (gotas sobre la membrana ePTFE) fue de
$8 \pm 2$\,min, tras los cuales la señal recuperó el 95\,\% de su
lectura de referencia.

% -----------------------------------------------------------------------------
% 3.3. COMPENSACIÓN Y VALIDACIÓN
% -----------------------------------------------------------------------------
\subsection{Curvas de compensación y validación contra patrón}
\label{subsec:cap2_compensacion_validacion}

La aplicación del modelo de compensación multivariado
(Ecuación~\ref{eq:compensacion_ndir}) con los coeficientes ajustados
durante las corridas experimentales ($\beta = 2.1 \times 10^{-4}$,
$\gamma = -0.45$\,ppm\,K$^{-1}$) redujo sistemáticamente la
desviación respecto al método de referencia GC-TCD
(Tabla~\ref{tab:validacion_sensor}).

\begin{table}[htbp]
\centering
\caption{Indicadores de desempeño metrológico del nodo IoT tras
compensación cruzada, validados contra cromatografía de gases
(GC-TCD).}
\label{tab:validacion_sensor}
\begin{tabular}{lcc}
\toprule
\textbf{Métrica} & \textbf{\ce{CO2}} & \textbf{\ce{CH4}} \\
\midrule
Rango de medición (ppm) & 400--6000 & 0--5000 \\
N pares (sensor, referencia) & 36 & 36 \\
Pearson $R$ & 0.94 & 0.89 \\
MAPE (\%) & 8.3 & 11.7 \\
RMSE (ppm) & 312 & 287 \\
Incertidumbre expandida $U$ ($k=2$) (ppm) & 520 & 410 \\
\bottomrule
\end{tabular}
\end{table}

La dispersión de los pares de validación mostró una relación lineal
sustancial para \ce{CO2} (Figura~\ref{fig:validacion_co2}), con los
puntos distribuidos uniformemente a lo largo de la recta 1:1 y sin
tendencia sistemática en los residuos estandarizados. Para \ce{CH4},
la correlación fue ligeramente inferior, atribuible a la menor
resolución del sensor NDIR en el rango bajo ($<< 500$\,ppm) y a la
presencia de interferencias por compuestos orgánicos volátiles
(co-elución espectral en la banda de 3.3\,$\mu$m).

\begin{figure}[htbp]
\centering
% \includegraphics[width=0.45\textwidth]{figuras/validacion_co2}
% \includegraphics[width=0.45\textwidth]{figuras/validacion_ch4}
\caption{Validación del nodo IoT contra GC-TCD. Izquierda: dispersión
de \ce{CO2} ($R = 0.94$, MAPE $= 8.3$\,\%). Derecha: dispersión de
\ce{CH4} ($R = 0.89$, MAPE $= 11.7$\,\%). Línea continua: ajuste
lineal; línea punteada: recta 1:1.}
\label{fig:validacion_co2}
\end{figure}

El MAPE obtenido para \ce{CO2} (8.3\,\%) y \ce{CH4} (11.7\,\%) cumple
con el umbral de aceptación establecido en la hipótesis H1
(MAPE $\leq 15$\,\%), aunque el valor de \ce{CH4} se aproxima al límite
superior, sugiriendo que sensores NDIR de mayor selectividad
espectral o algoritmos de corrección por interferentes podrían
mejorar la precisión en futuras iteraciones del hardware.

% -----------------------------------------------------------------------------
% 3.4. DATASET EXPERIMENTAL
% -----------------------------------------------------------------------------
\subsection{Dataset experimental: estadísticas descriptivas}
\label{subsec:cap2_dataset}

El despliegue del nodo IoT durante los seis ciclos experimentales
generó un dataset de series temporales de alta frecuencia (30\,s de
resolución), con una continuidad operativa del 97.4\,\% y una tasa de
pérdida de paquetes MQTT del 2.1\,\%. La Tabla~\ref{tab:estadisticas_dataset}
resume las estadísticas descriptivas de las concentraciones registradas
por dieta y gas.

\begin{table}[htbp]
\centering
\caption{Estadísticas descriptivas del dataset experimental de
emisiones. Valores en ppm.}
\label{tab:estadisticas_dataset}
\begin{tabular}{lcccccc}
\toprule
\textbf{Gas} & \textbf{Dieta} & \textbf{$N$} & \textbf{Media} &
\textbf{DE} & \textbf{Mín.} & \textbf{Máx.} \\
\midrule
\ce{CO2} & D1 & 2\,880 & 2\,340 & 1\,420 & 410 & 6\,000 \\
\ce{CO2} & D4 & 2\,880 & 1\,890 & 1\,150 & 400 & 5\,950 \\
\ce{CH4} & D1 & 2\,880 & 4\,560 & 6\,120 & 0 & 27\,500 \\
\ce{CH4} & D4 & 2\,880 & 2\,980 & 4\,870 & 0 & 19\,800 \\
\bottomrule
\end{tabular}
\end{table}

Las series temporales de \ce{CO2} exhibieron un patrón de acumulación
monotónica durante las fases de cierre hermético, con tasas de cambio
instantáneas que oscilaron entre 150 y 625\,ppm\,min$^{-1}$ en los
primeros 10 min de confinamiento para la dieta D1. Las concentraciones
de \ce{CH4}, por su parte, mostraron un comportamiento de umbral
característico: permanecieron en niveles basales ($<< 200$\,ppm) durante
un periodo de latencia variable (3--12 min), tras el cual se activó una
fase de producción exponencial que alcanzó concentraciones de hasta
27\,500\,ppm en el experimento con dieta D1 (Figura~\ref{fig:series_temporales}).

\begin{figure}[htbp]
\centering
% \includegraphics[width=0.8\textwidth]{figuras/series_temporales_d1_d4}
\caption{Series temporales representativas de \ce{CO2} (línea continua)
y \ce{CH4} (línea punteada) durante ciclos de cierre hermético.
Panel superior: dieta D1. Panel inferior: dieta D4. Nota: los ejes
de \ce{CH4} están escalados por 0.1 para facilitar la comparación
visual.}
\label{fig:series_temporales}
\end{figure}

La diferencia en el comportamiento de \ce{CH4} entre dietas fue
estadísticamente significativa: la dieta D1 alcanzó picos de
metano 1.4 veces mayores que la D4 ($p < 0.01$, prueba $t$ de Student),
lo cual se atribuye a la mayor densidad de carga biodegradable y la
formación más rápida de núcleos anaerobios en el sustrato proteico
de alta energía. Este hallazgo valida empíricamente la hipótesis de
que el \ce{CH4} actúa como indicador de condiciones operativas
subóptimas, particularmente en sustratos de alta carga metabólica.

